% BHU PYQ -- Manually transcribed & error-corrected
\documentclass[11pt,a4paper]{article}

\usepackage[a4paper,top=2.5cm,bottom=2.5cm,left=2.8cm,right=2.8cm,headheight=14pt]{geometry}
\usepackage{amsmath,amssymb}
\usepackage[T1]{fontenc}
\usepackage[utf8]{inputenc}
\usepackage{lmodern}
\usepackage{microtype}
\usepackage{fancyhdr}
\usepackage{enumitem}
\usepackage{hyperref}
\usepackage{lastpage}
\usepackage{parskip}

\hypersetup{colorlinks=true,urlcolor=black,linkcolor=black,
  pdftitle={BPT-301: Optics -- BHU PYQ},pdfauthor={Banaras Hindu University}}

\pagestyle{fancy}\fancyhf{}
\renewcommand{\headrulewidth}{0.4pt}\renewcommand{\footrulewidth}{0.4pt}
\fancyhead[L]{\small   \textit{Banaras Hindu University}}
\fancyhead[R]{\small   \textit{BPT-301: Optics}}
\fancyfoot[C]{\small Page  \thepage\ of \pageref{LastPage}}


\newlist{parts}{enumerate}{2}
\setlist[parts,1]{label=(\alph*),leftmargin=*,itemsep=5pt,topsep=3pt}
\setlist[parts,2]{label=(\roman*),leftmargin=*,itemsep=3pt,topsep=2pt}

\newcommand{\pts}[1]{\hfill   \textit{[#1]}}


\begin{document}


\begin{center}
  {\large\bfseries BANARAS HINDU UNIVERSITY}\\[2pt]
  {\normalsize B.Sc. (Hons.) Semester III Examination, 2022-23}\\[6pt]
  \rule{0.8\linewidth}{0.5pt}\\[6pt]
  {\large\bfseries Paper No. BPT-301: Optics}\\[4pt]
  {\normalsize Time: 3 Hours\hfill Maximum Marks: 70}
\end{center}
\rule{\linewidth}{0.4pt}
\smallskip



\noindent   \textit{Instructions: Answer   \textbf{five} questions including Question~1 which is compulsory. Symbols have their usual meanings.}
\bigskip




\noindent   \textbf{Question 1.} Answer any   \textbf{five} of the following: \pts{5 $ \times$ 4 = 20}

\begin{parts}
  \item An interference pattern from two point sources ($\lambda = 5000\, \text{\AA}$) separated by a distance $d = 0.05\, \text{mm}$ is produced on a screen $PP'$ at a distance $D = 20\, \text{cm}$ from source $S_1$. The plane of the screen $PP'$ is perpendicular to the line joining $S_1$ and $S_2$. What will be the shape of the fringes? What will be the order and the nature of the central fringe?
  \item A left circularly polarized beam ($\lambda = 5893\, \text{\AA}$) is incident normally on a calcite crystal ($n_o = 1.65836$ and $n_e = 1.48641$) of thickness $0.005141\, \text{mm}$ with its optic axis parallel to the surface. What will be the state of polarization of the emergent beam?
  \item Discuss the state of polarization when the $x$- and $y$-components of the electric field are given by the following equations:
    \[ E_x = E_0 \sin(kz - \omega t + \pi/4), \quad E_y = E_0 \sin(kz - \omega t - \pi/4) \]
  \item In a double-slit diffraction experiment, consider a slit width $b = 5  \times 10^{-2}\, \text{cm}$ and a separation between identical points of the two slits $d = 0.1\, \text{cm}$, illuminated by a monochromatic light of wavelength $6  \times 10^{-5}\, \text{cm}$. How many interference minima will occur within the central diffraction maximum?
  \item A straight edge is kept midway between a monochromatic source ($\lambda = 5000\, \text{\AA}$) and a screen. The distance between the source and the screen is $50\, \text{cm}$. Calculate the distance of the first two maxima and the first two minima from the edge of the geometrical shadow.
  \item Calculate the minimum number of lines in a grating that will resolve the sodium doublet ($\lambda_1 = 5890\, \text{\AA}$ and $\lambda_2 = 5896\, \text{\AA}$) in the second order.
  \item Using sodium light of wavelength $\lambda = 5893\, \text{\AA}$, interference fringes are formed by reflection from a thin air wedge-shaped film. Ten fringes are observed in a distance of $1\, \text{cm}$ when viewed perpendicularly. Calculate the angle of the wedge.
\end{parts}

\medskip\rule{\linewidth}{0.2pt}




\noindent   \textbf{Question 2.}

\begin{parts}
  \item In the figure of Lloyd's mirror given below, $S$ is a point source emitting waves of frequency $6  \times 10^{14}\, \text{Hz}$, $AB$ is the mirror and $LOM$ is the screen. The distances $SP$, $PA$, $AB$ and $BO$ are $1\, \text{mm}$, $5\, \text{cm}$, $5\, \text{cm}$ and $190\, \text{cm}$ respectively. Determine the position of the region where the fringes will appear and the number of fringes in the region. \pts{6}
  \item Discuss the working of a Michelson interferometer and show that the highest order fringe appears at the centre. How can the difference between two close wavelengths of a bichromatic source be calculated using a Michelson interferometer? \pts{8}
\end{parts}

\medskip\rule{\linewidth}{0.2pt}




\noindent   \textbf{Question 3.}

\begin{parts}
  \item Discuss the intensity distribution of Fraunhofer diffraction due to $N$ slits. Deduce the expression for the width of the principal maxima. How many minima and secondary maxima are present between two closest principal maxima? \pts{8}
  \item In the Newton's ring arrangement, if the incident light consists of two wavelengths $4000\, \text{\AA}$ and $4002\, \text{\AA}$, calculate the distance from the point of contact at which the rings will disappear. The radius of curvature of the lens surface is $400\, \text{cm}$. \pts{6}
\end{parts}
\hfill   \textbf{OR}

\begin{parts}
  \item Derive the expression and discuss the intensity distribution of a Fabry-P\'erot interferometer. Discuss how the sharpness of the fringes depends on the reflectivity of the mirrors. Give a brief comment on the resolving power of a Fabry-P\'erot interferometer. \pts{8}
  \item Discuss the Rayleigh criterion of resolving power. Find the resolving power of a microscope. \pts{6}
\end{parts}

\medskip\rule{\linewidth}{0.2pt}




\noindent   \textbf{Question 4.}

\begin{parts}
  \item Explain the working of a Retardation Plate, a Quarter-Wave Plate (QWP), and a Half-Wave Plate (HWP). \pts{6}
  \item Discuss the construction and working of a Nicol prism. \pts{4}
  \item What is optical activity? Explain Fresnel's theory of optical rotation. How is specific rotation experimentally determined using a Laurent half-shade polarimeter? \pts{4}
\end{parts}
\hfill   \textbf{OR}

\begin{parts}
  \item The plane of polarization of a plane-polarized light is turned through an angle $12.6^\circ$ passing through a $10\%$ sugar solution of length $20\, \text{cm}$. Calculate the specific rotation. \pts{6}
  \item What are Einstein coefficients? Derive the relation between them. \pts{8}
\end{{parts}}

\medskip\rule{\linewidth}{0.2pt}




\noindent   \textbf{Question 5.}

\begin{parts}
  \item Discuss the intensity distribution due to a single-slit diffraction and find the width of the central maximum. \pts{14}
\end{parts}

\vfill
\rule{\linewidth}{0.4pt}

\begin{center}\small
  \href{https://bhu-exam.vercel.app/}{Click here to visit our website}\\[4pt]
  \href{https://me-aryan.vercel.app/}{Click here to visit our contributor}\\[4pt]
  \href{https://quashan-me.vercel.app/}{Click here to visit our contributor}
\end{center}

\end{document}
